\documentclass[11pt,a4paper]{article}
\usepackage[T1]{fontenc}
\usepackage[utf8]{inputenc}
\usepackage{amsmath,amssymb,xcolor}
\usepackage[margin=25mm]{geometry}
\usepackage[colorlinks=true,urlcolor=blue]{hyperref}
\definecolor{nexorablue}{HTML}{2459B5}
\setlength{\parindent}{0pt}
\setlength{\parskip}{7pt}
\setlength{\emergencystretch}{2em}
\renewcommand{\familydefault}{\sfdefault}
\begin{document}

{\small\bfseries\color{nexorablue} NEXORA CPGE}\par
{\small MPSI / NEX-CIR-03 / Énoncé}\par
{\LARGE\bfseries Une sortie avec et sans charge}\par
\vspace{0.25em}\hrule\vspace{0.6em}

Une source impose $V_A=12\,\mathrm V$ par rapport à la masse. Une résistance de $2\,\mathrm{k\Omega}$ relie $A$ à $B$, et une de $4\,\mathrm{k\Omega}$ relie $B$ à la masse. Calculer $V_B$. On ajoute ensuite une charge de $4\,\mathrm{k\Omega}$ entre $B$ et la masse. Recalculer $V_B$ et les trois courants dans les résistances.


\end{document}
